\chapter{Introducción}

\section{Contexto General y Planteamiento del Problema}

La electroencefalografía (EEG) ha sido, durante décadas, una de las herramientas fundamentales de la neurociencia clínica y experimental. Su alta resolución temporal —en el orden de los milisegundos—, su naturaleza no invasiva y su costo relativamente bajo frente a técnicas como la resonancia magnética funcional o la tomografía por emisión de positrones la han consolidado como el estándar para estudiar la dinámica cerebral en tiempo real \cite{buzsaki2006rhythms}. Sus aplicaciones abarcan el diagnóstico de epilepsia y trastornos del sueño, la investigación en neurociencia cognitiva y el desarrollo de interfaces cerebro-computador para comunicación asistida y neurorrehabilitación.

Ahora bien, la adquisición de EEG es un proceso frágil. En la práctica clínica y de investigación, la pérdida parcial o total de información en uno o más electrodos es frecuente. Estas fallas pueden deberse al aumento de la impedancia de contacto por transpiración o secado del gel conductor, a fallas en el amplificador como la saturación de voltaje, a movimientos bruscos del paciente que desconectan los sensores, o simplemente al descarte deliberado de canales que presentan una relación señal-ruido inaceptablemente baja durante la etapa de preprocesamiento \cite{nunez2006electric}.

La pérdida de canales no es un problema cosmético: compromete la integridad de todo análisis posterior. La ausencia de información espacial distorsiona la topografía de la señal, afecta la localización de fuentes neuronales y altera biomarcadores electrofisiológicos como los Potenciales Relacionados a Eventos (ERP) y la Densidad Espectral de Potencia (PSD) \cite{luck2014introduction}. En sistemas modernos de aprendizaje automático aplicados a interfaces cerebro-computador, la pérdida impredecible de canales cambia la dimensionalidad de entrada de los clasificadores y obliga, con frecuencia, a descartar ensayos completos de datos clínicamente valiosos \cite{lotte2018review}. La interpolación precisa de estos canales es, por tanto, un paso ineludible del flujo de preprocesamiento, y su calidad condiciona directamente la validez de cualquier conclusión neurofisiológica o clínica posterior.

\section{Limitaciones de los Enfoques Tradicionales}

Históricamente, el problema de la interpolación de canales EEG ha sido abordado mediante métodos espaciales o estadísticos. El enfoque más extendido y utilizado por defecto en bibliotecas de análisis como EEGLAB \cite{delorme2004eeglab} o MNE-Python \cite{gramfort2013meg,gramfort2014mne} es la interpolación por \textit{spherical splines} \cite{perrin1989}. Este método, junto con otros enfoques basados en distancia euclidiana como la interpolación por vecinos más cercanos o las funciones de base radial, asume que la cabeza humana es una esfera conductora perfecta y estima el voltaje del canal faltante proyectando geométricamente los voltajes de los electrodos circundantes sobre dicha superficie.

Aunque estos métodos son rápidos y matemáticamente elegantes, adolecen de una limitación fundamental: operan estrictamente en el dominio espacial estático para cada instante de tiempo individual. Al hacerlo, ignoran dos realidades neurofisiológicas esenciales. En primer lugar, la conectividad funcional real del cerebro no es isotrópica: la actividad eléctrica sigue rutas anatómicas complejas que no siempre se correlacionan con la distancia física en el cuero cabelludo, de modo que un modelo esférico resulta una simplificación excesiva \cite{sporns2005human}. En segundo lugar, las señales EEG son generadas por osciladores neuronales sujetos a leyes físicas de continuidad y momento. Al reconstruir un canal muestra por muestra de forma aislada, los métodos puramente espaciales son incapaces de preservar la dinámica de alta frecuencia, generando trazos con ruido de fase y distorsiones en los picos transitorios \cite{lai2005controlled}.

Los métodos estadísticos de descomposición —Análisis de Componentes Independientes (ICA) y Análisis de Componentes Principales (PCA)— proyectan las señales observadas sobre fuentes latentes y reconstruyen el canal perdido a partir de esa proyección \cite{jung2001imaging}. Sin embargo, cuando varios canales contiguos fallan simultáneamente, estos métodos tienden a inyectar ruido global o a perder los picos transitorios locales que definen la morfología fina de los ERP.

\section{El Advenimiento del Procesamiento de Señales en Grafos}

Para superar estas limitaciones, en los últimos años ha emergido con fuerza el marco del Procesamiento de Señales en Grafos (GSP, por sus siglas en inglés) \cite{shuman2013,ortega2018}. El GSP ofrece una perspectiva natural y poderosa para lidiar con datos que residen en dominios irregulares no euclidianos. En lugar de modelar la cabeza como un espacio continuo tridimensional, GSP modela la red de electrodos como un grafo discreto $\mathcal{G} = (\mathcal{V}, \mathcal{E}, \mathbf{W})$, donde cada electrodo es un nodo y las relaciones de similitud o distancia fisiológica entre ellos se codifican en las aristas mediante núcleos de distancia gaussianos o basados en correlación.

Bajo este paradigma, el voltaje medido en un instante se interpreta como una señal de grafo, y conceptos clásicos como la Transformada de Fourier, el filtrado paso-bajo y el muestreo se trasladan al dominio del grafo mediante el operador Laplaciano $\mathbf{L}$. La intuición es simple: las señales en nodos fuertemente conectados tienden a ser similares. Al minimizar una función de costo que penaliza las variaciones abruptas sobre la topología del grafo, los datos faltantes se reconstruyen con mayor fidelidad que usando solo la distancia euclidiana \cite{narang2013}.

Una extensión aún más avanzada es la regularización espacio-temporal. Trabajos recientes han introducido enfoques como la Regularización Temporal de Suavidad de Sobolev (TRSS, por sus siglas en inglés) \cite{giraldo2022}, que optimiza conjuntamente la topología espacial y la derivada temporal de la señal. Al penalizar saltos bruscos en el tiempo de forma simultánea a los saltos en el espacio, TRSS impone restricciones fisiológicas reales: las neuronas no cambian su potencial de membrana instantáneamente, sino que obedecen dinámicas de integración continua. Esta restricción dual —espacial y temporal— es lo que la teoría predice como el factor diferencial frente a los métodos puramente espaciales.

\subsection{Mención del Paradigma de Aprendizaje Profundo}
\label{sec:ia_mencion}

Existe además un cuerpo creciente de literatura que aborda la interpolación de señales EEG mediante redes neuronales profundas, incluyendo arquitecturas como autoencoders convolucionales \cite{banville2021uncovering}, redes generativas antagónicas para imputación de series temporales \cite{yoon2019time}, y más recientemente Redes Neuronales de Grafos que operan directamente sobre la topología de electrodos \cite{defferrard2016}. Estos enfoques han mostrado resultados prometedores en escenarios controlados, beneficiándose de la capacidad de las redes profundas para aprender representaciones latentes complejas a partir de grandes volúmenes de datos.

No obstante, estos métodos presentan limitaciones que motivan su exclusión deliberada del alcance de esta tesis. En primer lugar, son inherentemente dependientes de los datos de entrenamiento, lo que introduce riesgos de falta de generalización al transferir modelos entre poblaciones, equipos o paradigmas experimentales distintos. En segundo lugar, carecen de garantías formales de convergencia y de interpretabilidad sobre el mecanismo de interpolación, una propiedad crítica en contextos clínicos donde cada decisión debe ser trazable y justificable. En tercer lugar, su costo computacional de entrenamiento es órdenes de magnitud superior al de los métodos algorítmicos aquí considerados. Por estas razones, la presente investigación se circunscribe exclusivamente a métodos algorítmicos con fundamento matemático explícito —geométricos, estadísticos y basados en GSP— cuya operación no requiere entrenamiento y cuyas propiedades de convergencia y regularización son analíticamente caracterizables.

\section{La Brecha en la Literatura y la Motivación del Trabajo}

A pesar del potencial teórico de los algoritmos GSP y TRSS, su validación empírica en la literatura de señales biomédicas presenta limitaciones metodológicas relevantes. Se identifican tres brechas que motivan el diseño experimental de esta tesis.

El primer problema es que los estudios que proponen nuevos algoritmos GSP suelen compararse contra \textit{spherical splines} usando las configuraciones por defecto de los paquetes de software, sin realizar una búsqueda rigurosa de hiperparámetros para los métodos de referencia. Esto puede inflar artificialmente las ventajas del método propuesto y distorsionar la percepción del verdadero estado del arte \cite{saldivia2026benchmark}. El segundo problema es que la mayoría de los trabajos restringe su evaluación a métricas de error de amplitud global —como el RMSE o el MAE— y aplica escenarios de desconexión simples, tales como la remoción de canales aleatorios aislados en bases de datos pequeñas. Rara vez se investiga el impacto morfológico de la reconstrucción mediante métricas temporales como el \textit{Dynamic Time Warping} o espectrales como la Distancia Espectral Logarítmica. El tercer problema es la escasez de comparaciones contra implementaciones comunitarias maduras. En particular, la interpolación de MNE-Python incorpora más de una década de refinamientos de ingeniería —normalización de coordenadas, regularización diagonal adaptativa, aumentación del sistema con restricciones lineales, pseudoinversa con tolerancia adaptativa y un mecanismo de reintento automático— que la hacen sustancialmente más robusta que una implementación ingenua de \textit{spherical splines}. Ignorar esta implementación equivale a comparar contra un método de referencia debilitado artificialmente.

La motivación concreta de este trabajo nace directamente de estas brechas. El objetivo no es diseñar desde cero un nuevo marco teórico matemático, sino someter a la frontera del conocimiento actual —métodos geométricos, GSP y regularización temporal— a la evaluación empírica más rigurosa, exhaustiva y justa posible. Al utilizar el sistema moderno de optimización bayesiana Optuna \cite{akiba2019optuna} para llevar todos los algoritmos a su máximo límite empírico bajo protocolos evaluados en múltiples conjuntos de datos, esta tesis busca revelar el verdadero punto de cruce entre el paradigma clásico y el paradigma de grafos.

\begin{figure}[htbp]
\centering
\includegraphics[width=0.95\textwidth]{ch1_thesis_roadmap}
\caption[Mapa lógico de la tesis]{Mapa lógico de la tesis. La contribución se organiza como una cadena de ingeniería: caracterizar el problema de canales faltantes, comparar métodos espaciales y espacio--temporales bajo un protocolo común, y transformar los resultados en una frontera práctica de decisión entre MNE-Python y TRSS.}
\label{fig:ch1_thesis_roadmap}
\end{figure}

\subsection{La pregunta central: ¿por qué funciona tan bien la interpolación de MNE-Python?}
\label{sec:why_mne_works}

Una pregunta recurrente que emerge del análisis preliminar —y que motiva una sección dedicada en el marco teórico de esta tesis— es por qué la función de interpolación de MNE-Python produce resultados tan competitivos, incluso frente a métodos teóricamente superiores como TRSS. La respuesta, que se desarrolla en profundidad en el Capítulo~2, reside en que la implementación de MNE no es una simple transcripción del artículo original de Perrin y colaboradores \cite{perrin1989}, sino que incorpora un conjunto de mejoras de ingeniería acumuladas durante más de diez años de desarrollo como código abierto.

Estas mejoras —normalización de coordenadas a esfera unitaria, uso global de canales, regularización diagonal adaptativa, aumentación con restricciones lineales, pseudoinversa con tolerancia adaptativa y reintento automático— no aparecen descritas en la literatura académica, pero están presentes en el código fuente de MNE. Son las responsables de que un método formulado en 1989 siga siendo hoy una referencia competitiva. Esta tesis dedica una sección completa del marco teórico (Capítulo~2) a documentar estos mecanismos, y una validación confirmatoria independiente (Capítulo~6) a cuantificar la brecha real entre TRSS y MNE.

\section{Cronología del Proyecto}
\label{sec:cronologia}

La investigación que sustenta esta tesis no siguió un camino lineal, sino un proceso iterativo de descubrimiento guiado por los resultados experimentales que se iban obteniendo. Documentar esta trayectoria es esencial para comprender las decisiones metodológicas que dieron forma al protocolo final, y se presenta a continuación de manera cronológica.

\subsection{Fase 1 — Selección de Métodos y Revisión del Estado del Arte}

Durante los primeros dos meses se realizó una revisión sistemática de la literatura sobre interpolación de señales en dominios irregulares. Este proceso permitió identificar tres familias metodológicas bien diferenciadas: métodos geométricos o espaciales (\textit{spherical splines}, vecinos más cercanos, funciones de base radial), métodos estadísticos de descomposición (ICA, PCA, interpolación lineal) y métodos basados en GSP (regularización de Tikhonov en grafos, regularización de Sobolev, TRSS). De cada familia se seleccionaron los representantes más citados y con fundamento matemático más sólido, resultando en un total de dieciocho métodos candidatos para implementación y evaluación comparativa.

\subsection{Fase 2 — Implementación del Sistema de Experimentación}

Los siguientes dos meses se dedicaron al desarrollo de un sistema modular en Python que unificara la interfaz de todos los métodos bajo una interfaz de programación común. Cada método recibía una matriz de datos $\mathbf{X} \in \mathbb{R}^{C \times T}$ —con $C$ canales y $T$ muestras temporales— junto con una máscara binaria que indicaba qué canales debían tratarse como faltantes, y entregaba como resultado la matriz reconstruida. En paralelo se implementaron los constructores de grafos necesarios —incluyendo k-NN con pesos gaussianos, grafos de correlación funcional y grafos de distancia completa—, los núcleos de interpolación GSP y el resolvedor TRSS con regularización dual mediante los parámetros $\alpha$ y $\beta$. Esta fase fue particularmente demandante porque requirió validar que cada implementación propia reprodujera exactamente el comportamiento descrito en su artículo original antes de incorporarla al sistema comparativo.

\subsection{Fase 3 — Barrido Exploratorio Inicial}

Con el sistema funcional, se ejecutó un barrido exploratorio sobre escenarios de pérdida aleatoria al 20\% en los tres conjuntos de datos. Los resultados revelaron un hallazgo que redefinió el curso de la investigación: TRSS mostraba ventajas del orden de 5 a 15~dB en SNR sobre las implementaciones propias de \textit{spherical splines}, pero esa brecha se reducía de manera importante al comparar contra la implementación comunitaria de referencia de MNE-Python, que incorpora una década de refinamientos de ingeniería. Este hallazgo fue el disparador de dos decisiones metodológicas cruciales: incluir a MNE como referencia confirmatoria independiente en una fase separada del experimento, y diseñar el resto del protocolo para distinguir nítidamente entre resultados exploratorios y confirmatorios.

\subsection{Fase 4 — Selección Fina y Diseño del Espacio de Búsqueda}

Los resultados del barrido permitieron descartar con confianza los métodos claramente inferiores —como la interpolación lineal o el vecino más cercano sin ponderación— y concentrar los recursos de optimización en los métodos con desempeño prometedor. Para cada familia se definió un espacio de búsqueda de hiperparámetros diferenciado: la escala del núcleo gaussiano $\sigma$ y el número de vecinos $k$ para los métodos de grafo, los coeficientes de regularización espacial $\alpha$ y temporal $\beta$ para TRSS, y el orden del spline $m$ junto con la regularización $\lambda$ para los \textit{spherical splines}.

\subsection{Fase 5 — Optimización Bayesiana con Optuna}

Durante los siguientes dos meses se integró Optuna con muestreo TPE (\textit{Tree-structured Parzen Estimator}) como motor de optimización \cite{akiba2019optuna}. Se ejecutaron treinta intentos independientes por cada método de grafo, diez intentos con el algoritmo multiobjetivo NSGA-II para los \textit{spherical splines}, y barridos de grilla fina para TRSS dado su mayor costo computacional. En total se evaluaron más de quince mil configuraciones paramétricas distintas. Esta fase constituye el núcleo experimental de la tesis: sus resultados son los que se reportan en el Capítulo~6.

\subsection{Fase 6 — Validación Confirmatoria contra MNE-Python}

En una fase independiente —diseñada para no contaminar los resultados exploratorios— se ejecutó una validación donde TRSS, con sus hiperparámetros congelados tras la optimización de la fase anterior, se comparó directamente contra la implementación de referencia de MNE-Python utilizando los parámetros por defecto que cualquier usuario de MNE obtendría sin ajuste adicional. Esta validación se estratificó por conjunto de datos, modo de pérdida y severidad, y empleó pruebas estadísticas pareadas —Wilcoxon signed-rank con corrección de Benjamini-Hochberg para comparaciones múltiples— con el fin de establecer significancia estadística de manera rigurosa.

\subsection{Fase 7 — Ablación y Documentación}

La fase final consistió en un estudio de ablación para cuantificar la contribución relativa de cada componente de TRSS —el término espacial $\alpha$ frente al término temporal $\beta$—, la generación sistemática de todas las figuras y tablas que aparecen en este documento, y la redacción de la tesis propiamente dicha.

\section{Objetivos}

\subsection{Objetivo General}

Comparar el rendimiento espacial, temporal y espectral de los métodos de reconstrucción de canales EEG basados en GSP con regularización temporal frente a las referencias geométricas —incluida MNE-Python— bajo un marco experimental riguroso, reproducible y optimizado mediante optimización bayesiana.

\subsection{Objetivos Específicos}

\begin{itemize}
    \item \textbf{OE1 — Sistema unificado de evaluación.} Implementar dieciocho métodos de interpolación en un sistema modular con interfaz común, definir protocolos de pérdida aleatoria y contigua (10\%--40\%), e integrar Optuna para la optimización imparcial de todos los algoritmos (más de 15\,000 configuraciones exploradas).

    \item \textbf{OE2 — Evaluación multi-métrica y fisiológica.} Cuantificar el desempeño mediante seis métricas primarias (MAE, RMSE, SNR, DTW, LSD, Coherencia) sobre tres conjuntos de datos y 120 escenarios de pérdida, y analizar la preservación morfológica y espectral mediante métricas auxiliares, visualizaciones temporales y potencia por banda.

    \item \textbf{OE3 — Trazabilidad y reproducibilidad.} Generar un repositorio de artefactos computacionales —código, configuraciones, registros de optimización y resultados— que garantice la trazabilidad completa de cada afirmación cuantitativa hasta el archivo que la origina.
\end{itemize}

\section{Hipótesis de Trabajo}

\begin{quote}
\textbf{H1 — Hipótesis principal.} La regularización espacio-temporal (TRSS) permite reconstruir canales EEG faltantes con mayor precisión en métricas de amplitud y ajuste temporal que los métodos geométricos, especialmente bajo pérdida contigua severa.
\end{quote}

\begin{quote}
\textbf{H2 — Hipótesis de ventaja práctica de MNE.} En regímenes favorables —montajes densos ($\geq 60$ canales), pérdida aleatoria baja o señales espacialmente suaves—, MNE-Python iguala o supera a TRSS en preservación espectral y tiempo de cómputo, con diferencias de reconstrucción operacionalmente pequeñas.
\end{quote}

\section{Alcance y Aportes de la Investigación}

El alcance de este trabajo abarca la ejecución computacional a gran escala del sistema de evaluación sobre tres conjuntos de datos, ciento veinte escenarios de falla y dieciocho métodos de interpolación. A diferencia de las evaluaciones acotadas que predominan en la literatura, aquí se consolidan miles de ensayos EEG bajo más de quince mil configuraciones paramétricas exploradas mediante Optuna.

Para garantizar la reproducibilidad, la ejecución se ancla en un registro exhaustivo de corridas experimentales que fija explícitamente las combinaciones de conjuntos de datos, modos de pérdida, semillas, tipos de grafo y familias de métodos. Este diseño evita las evaluaciones improvisadas y permite trazar cada afirmación cuantitativa reportada en los capítulos posteriores hasta el artefacto experimental que la origina, según se detalla en el Apéndice~D.

Los aportes principales se articulan en cuatro ejes. Primero, se proporciona evidencia estadística de que TRSS mejora de forma robusta las métricas de amplitud y ajuste temporal frente a referencias geométricas en regímenes difíciles, una vez que todos los métodos han sido optimizados. Segundo, se muestra que bajo pérdida contigua severa (30\% o más de los canales) la restricción temporal puede compensar la pérdida de información espacial local. Tercero, se documentan en formato académico las seis mejoras de ingeniería que convierten a MNE-Python en una referencia inesperadamente competitiva. Cuarto, se entrega un sistema unificado en Python, rastreable y reproducible para ejecutar interpolaciones GSP con optimización bayesiana automatizada sobre datos EEG.

\section{Estructura de la Tesis}

Para desarrollar el análisis propuesto, el resto del documento se organiza de la siguiente manera. El Capítulo~2 presenta los fundamentos neurofisiológicos del EEG, los pilares matemáticos del Procesamiento de Señales en Grafos —incluyendo el Laplaciano, la Transformada de Fourier en Grafos y la construcción de grafos—, el estado del arte en interpolación de señales biomédicas, y una sección dedicada al análisis detallado del mecanismo interno de la función de interpolación de MNE-Python. El Capítulo~3 detalla el protocolo experimental: los tres conjuntos de datos empleados, el modelado matemático de los dieciocho algoritmos evaluados, los métodos de simulación de daños y la arquitectura del motor de optimización bayesiana. El Capítulo~4 documenta la arquitectura de software del sistema experimental, la trazabilidad entre módulos y la gestión de artefactos. El Capítulo~5 formaliza el diseño experimental: variables independientes, dependientes y controladas, el protocolo de validación y el plan de análisis estadístico. El Capítulo~6 expone los hallazgos empíricos, incluyendo comparativas globales de error, curvas de degradación por pérdida creciente, visualizaciones temporales representativas, comparación espectral mediante PSD, el estudio de ablación de TRSS y la validación confirmatoria contra MNE-Python. El Capítulo~7 ofrece una discusión crítica sobre las causas metodológicas de los resultados obtenidos, interpreta los mecanismos subyacentes y transparenta las limitaciones del estudio. Finalmente, el Capítulo~8 sintetiza las conclusiones, valora el cumplimiento de las hipótesis y sugiere vías de investigación futura.

Los apéndices complementan el documento con derivaciones matemáticas completas (Apéndice~A), tablas extensas de resultados (Apéndice~B), el protocolo de reproducibilidad (Apéndice~C), el inventario de artefactos computacionales (Apéndice~D) y una adaptación de las tablas del documento original de tesis para referencia histórica (Apéndice~E).

\section{Tesis, Contribución y Criterio de Lectura}

La tesis que organiza este documento sostiene que la reconstrucción de canales EEG faltantes debe evaluarse como un problema simultáneamente espacial, temporal y espectral. En consecuencia, una evaluación basada únicamente en métricas de error de amplitud como el MAE o el RMSE no es suficiente para decidir si una interpolación es útil en aplicaciones neurofisiológicas reales, porque estas métricas son insensibles a las distorsiones de fase y a la preservación de la estructura fina del espectro de frecuencias.

Esta investigación adopta una posición deliberadamente conservadora y transparente: los resultados favorables a TRSS se reportan solo cuando pueden trazarse a artefactos reproducibles; las limitaciones se explicitan con honestidad; y se reconoce que la implementación de MNE-Python constituye una referencia competitiva en sus regímenes de operación óptimos.

La contribución central puede sintetizarse en una sola oración: \textit{esta tesis establece una evaluación comparativa reproducible y justa para la reconstrucción de canales EEG, mostrando que la regularización temporal sobre grafos —en particular TRSS— mejora métricas de amplitud y ajuste temporal bajo escenarios de pérdida severa, incluso después de optimizar rigurosamente los métodos de referencia, y documenta los mecanismos de ingeniería que hacen de MNE-Python una referencia comunitaria difícil de superar en regímenes favorables.}